\chapter{Nouns and pronouns}
\label{chap:nouns}

\section{Gender and articles}

\subsection{Gender}

Nouns in Niklorsk have masculine, feminine, or neuter gender. Gender is largely a lexical property of each noun, and the definite article is its primary overt indicator. The nominative articles are \nk{en} (masculine), \nk{de} (feminine), \nk{eg} (neuter), and \nk{der} (plural). The plural article is used regardless of the noun's gender. The article precedes the noun or noun phrase it governs. Its separate dative forms are \nk{om}, \nk{der}, \nk{egn}, and \nk{dæ}, respectively; it is otherwise indeclinable. Niklorsk has no indefinite article.

The following tendencies associate gender with a noun's referent or its written ending.

\begin{enumerate}[label=(\arabic*), wide]
    \item Nouns referring to people take the gender of the person they refer to. If the gender is unknown, or when using the noun to the specific group of people referred to by the noun, the noun is taken to be neuter. Most nouns referring to people can be used as masculine, feminine, or neuter, depending on the context, without any change in form. Some nouns, however, have distinct gendered forms.
    \begin{itemize}
        \item \nikp{en hond}{the (boy) child}, \nikp{de hond}{the (girl) child}, \nikp{eg hond}{the child}
        \item \nikp{en/de gläst}{the (male/female) doctor}, \nikp{eg gläst}{the doctor}
    \end{itemize}
    But:
    \begin{itemize}
        \item \nikp{en bund}{the king}, \nikp{de træ}{the queen}
        \item \nikp{en tadd}{the father}, \nikp{de mama}{the mother}
        \item \nikp{en Papps}{the Pope}
    \end{itemize}
    
    \item Nouns ending in \nk{-ll}, \nk{-tt} and \nk{-pp} are usually feminine, e.g., \nikp{de skäll}{the fish}, \nikp{de fjeltt}{the party}, \nikp{de flikk}{the name}, \nikp{de skäpp}{the ship}, \nikp{de kapp}{the coat}.
    
    \item Nouns ending in consonant clusters starting with \nk{r} or \nk{l} (except long \nk{l}), \nk{s} and \nk{g} are generally masculine.
    
    \item Nouns ending in \nk{-a} are usually feminine.
    \item Nouns ending in \nk{-e} are usually neuter.
\end{enumerate}

\subsection{Strong and weak articles}
\label{sec:article-reduction}

Definite articles have strong forms under contrastive emphasis and weak forms in ordinary unstressed noun phrases. Strong forms retain their full vowel qualities; weak forms attach prosodically to the following noun or adjective. In rapid speech, the nasal articles have syllabic consonant forms, as shown in Table~\ref{tab:weak-articles}. This attachment does not change their spelling or their gender, number, and case functions.

Weak articles have no independent lexical accent. A strong article pronounced as an independent stressed word is subject to the primary-stress and rhyme-licensing conditions in Chapter~\ref{chap:phonology}, \S\ref{sec:lexical-accent}.

\begin{table}[!ht]
    \caption{Weak and clitic forms of definite articles}
    \label{tab:weak-articles}
    \small
    \begin{tabularx}{\textwidth}{l X l X}
        \toprule
        \rowcolor{grammarTableHead}
        \textbf{Article} & \textbf{Function} & \textbf{Weak form} & \textbf{Rapid clitic form} \\
        \midrule
        \nk{en} & Masculine nominative & [\ipa{ən}] & [\ipa{n̩}] \\
        \nk{de} & Feminine nominative & [\ipa{də}] & [\ipa{d}] before a vowel; [\ipa{də}] elsewhere \\
        \nk{eg} & Neuter nominative & [\ipa{ə}] & [\ipa{ə}] \\
        \nk{der} & Plural nominative; feminine dative & [\ipa{dɐ}] & [\ipa{dɐ}] \\
        \nk{om} & Masculine dative & [\ipa{əm}] & [\ipa{m̩}] \\
        \nk{egn} & Neuter dative & [\ipa{əŋ}] & [\ipa{ŋ̩}] \\
        \nk{dæ} & Plural dative & [\ipa{dɐ}] & [\ipa{dɐ}] \\
        \bottomrule
    \end{tabularx}
\end{table}

Thus \nikp{en afte}{the book} has the rapid pronunciation [\ipa{n̩ ˈɑft}]. The article remains audible even though its vowel has no independent articulation. Weak \nk{der} and \nk{dæ} are homophonous, but their distribution in the noun phrase remains distinct. Strong forms are used for correction, contrast, and citation. The general prosodic conditions for weak forms are described in Chapter~\ref{chap:phonology}, \S\ref{sec:connected-speech}.

\subsection{Number}

Nouns are marked for number, with the unmarked form being singular. The plural is formed with one of the suffixes \nk{-(a)r}, \nk{-or}, \nk{-a}, \nk{-n}, or \nk{-er}, or by umlaut alone. A small class has a zero plural, with no change to the singular form. A restricted but productive class of loanwords takes the borrowed plural suffix \nk{-s}. Nouns in the \nk{-er} class combine the ending with umlaut of the stem vowel. The choice of plural class is lexical and must be memorized for each noun.

Plural endings preserve the position of the noun's lexical stress (see Chapter~\ref{chap:phonology}, \S\ref{sec:stress-prosody}). Thus \nk{afte} and \nk{aften} have initial stress, while \nk{magasin} and \nk{magasinnar} have stress on the third syllable. A change in vowel or consonant quantity during plural formation does not by itself move stress.

Ordinary productive plural formation preserves the stem's accent category unless a lexical paradigm specifies an alternation. The noun \nk{glus} has A1. Its plural \nk{glussar} retains that category in the paradigm and has surface $A_0$: its short vowel and geminate obstruent do not license the opposition. In \nk{magasin} and \nk{magasinnar}, lexical A1 is realised at the primary-stressed third syllable. Their sonorous stressed rhymes license the opposition, with a geminate sonorant in the plural. Stress preservation and accent licensing therefore have separate effects (Chapter~\ref{chap:phonology}, \S\ref{sec:lexical-accent-paradigms}).

Unstressed \nk{-ar} and \nk{-er} are both pronounced [\ipa{ɐ}], and plural \nk{-a} also retains [\ipa{ɐ}]. The classes remain distinct in their stem alternations and written endings. In \nk{glus} and \nk{glussar}, the plural has a shorter stem vowel and a long consonant as well as the ending [\ipa{ɐ}]. In \nk{bund} and \nk{bynder}, the changed stem vowel distinguishes the plural. Reduction of the ending does not remove these quantity and umlaut contrasts.

The plural \nk{-n} is silent. Singular \nikp{afte}{book} and plural \nikp{aften}{books} are both pronounced [\ipa{ˈɑft}] in ordinary speech; careful pronunciation retains final [\ipa{ə}] in both. Both have surface $A_0$ and remain homophonous because the short vowel and obstruent coda do not license a lexical accent contrast. Restoring the unstressed vowel does not change this restriction. The written \nk{-n} class remains distinct from the lexical zero-plural class described below. Number can also be indicated by articles, agreement, numerals, and context.

The following examples illustrate suffixal plurals without umlaut:

\begin{enumerate}
    \item Nouns taking \nk{-(a)r}: 
    
    Some nouns in this class have a shorter stem vowel and a long final consonant in the plural, as in \nikp{glus}{cat} $\rightarrow$ \nikp{glussar}{cats}

    \item Nouns taking \nk{-n}: \nikp{afte}{book} $\rightarrow$ \nikp{aften}{books}
\end{enumerate}

Nominalized participles always take the plural suffix \nk{-(a)r}. For example, \nikp{eg bördade}{the one eats} $\rightarrow$ \nikp{der bördader}{those who eat}, \nikp{eg hjaltagne}{the written thing} $\rightarrow$ \nikp{der hjaltagner}{the written things}.

Plural genitives follow the formation rules in \S\ref{sec:nominal-case}, e.g., \nikp{de blivt brug}{the man's house} $\rightarrow$ \nikp{de blivtes brug}{the men's house}.

The plural form is used to refer to more than one of something. Unlike in English, however, the plural form is not used in Niklorsk where the noun is preceded by a numeral or a quantifier. Thus we have \nikp{ty afte}{two books} and not \nk{ty aften}.

\subsubsection{Borrowed -s plurals}
\label{sec:plural-loan-s}

Recent loanwords are adapted to Niklorsk phonology and orthography and follow the ordinary rules of gender assignment. The borrowed plural suffix \nk{-s} is attached to the nativised stem. Thus English \emph{fan} is borrowed as \nk{fänn}, with the plural \nk{fänns}; retaining the plural ending does not preserve the English spelling.

Plurals in \nk{-s} frequently compete with native forms, particularly those in \nk{-(a)r}. The choice is lexical and may be conditioned by register, generation, or meaning. The plural ending does not determine the degree of phonological adaptation: a fully nativised stem can retain \nk{-s}, while a recent borrowing can take a native plural immediately. Table~\ref{tab:borrowed-s-plurals} gives examples of the competing forms.

\begin{table}[!ht]
    \caption{Borrowed \nk{-s} plurals and native alternatives}
    \label{tab:borrowed-s-plurals}
    \small
    \begin{tabularx}{\textwidth}{l l l >{\raggedright\arraybackslash}X}
        \toprule
        \rowcolor{grammarTableHead}
        \textbf{Singular} & \textbf{Plural in \nk{-s}} & \textbf{Native plural} & \textbf{Meaning} \\
        \midrule
        \nk{fänn} & \nk{fänns} & \nk{fännar} & fan, supporter \\
        \nk{äpp} & \nk{äpps} & \nk{äppar} & app \\
        \nk{klipp} & \nk{klipps} & \nk{klippar} & clip \\
        \nk{snäkk} & \nk{snäkks} & \nk{snäkkar} & snack \\
        \nk{spämm} & \nk{spämms} & \nk{spämmar} & spam \\
        \nk{strim} & \nk{strims} & \nk{strimmar} & stream \\
        \nk{läptåpp} & \nk{läptåpps} & \nk{läptåppar} & laptop \\
        \nk{skänn} & \nk{skänns} & \nk{skännar} & scan \\
        \nk{blåkk} & \nk{blåkks} & \nk{blåkkar} & blog \\
        \nk{slågan} & \nk{slågans} & \nk{slåganar} & slogan \\
        \bottomrule
    \end{tabularx}
\end{table}

The stem of \nk{strim} retains its long vowel in \nk{strims}. In the native plural \nk{strimmar}, the vowel is shortened and the following consonant is doubled, as in \nk{glus} $\rightarrow$ \nk{glussar}. These quantity changes do not move lexical stress. The recent loan \nk{strim} has default A1, which both plural forms preserve. The long vowel of \nk{strim} and \nk{strims} and the vowel--sonorant sequence of \nk{strimmar} all license its realisation, unlike the short vowel and geminate obstruent in \nk{glussar}.

The standard plural of \nikp{pådkasst}{podcast} is \nk{pådkasstar}; \nk{pådkassts} is a specialised variant. The native ending avoids the final consonant cluster of \nk{pådkassts}.

\subsubsection{The -er class with umlaut}
\label{sec:plural-er-umlaut}

Nouns in this class form the plural with \nk{-er} and umlaut of the stem vowel. Plurality is marked both by the changed stem vowel and by the ending. Table~\ref{tab:plural-er-umlaut} gives the forms and their stem alternations. In \nk{hus} $\rightarrow$ \nk{hysser}, the stem-final \nk{s} is also doubled.

\begin{table}[!ht]
    \caption{Plurals in \nk{-er} with umlaut}
    \label{tab:plural-er-umlaut}
    \small
    \begin{tabularx}{\textwidth}{l l >{\raggedright\arraybackslash}X >{\raggedright\arraybackslash}X}
        \toprule
        \rowcolor{grammarTableHead}
        \textbf{Singular} & \textbf{Plural} & \textbf{Meaning} & \textbf{Change} \\
        \midrule
        \nk{bund} & \nk{bynder} & king $\rightarrow$ kings & \nk{u} $\rightarrow$ \nk{y} \\
        \nk{hus} & \nk{hysser} & straw $\rightarrow$ straws & \nk{u} $\rightarrow$ \nk{y}, \nk{s} $\rightarrow$ \nk{ss} \\
        \nk{hjålt} & \nk{hjölter} & cake $\rightarrow$ cakes & \nk{å} $\rightarrow$ \nk{ö} \\
        \nk{hrådd} & \nk{hrödder} & group $\rightarrow$ groups & \nk{å} $\rightarrow$ \nk{ö} \\
        \nk{modd} & \nk{mydder} & meadow $\rightarrow$ meadows & \nk{o} $\rightarrow$ \nk{y} \\
        \nk{kapp} & \nk{käpper} & coat $\rightarrow$ coats & \nk{a} $\rightarrow$ \nk{ä} \\
        \bottomrule
    \end{tabularx}
\end{table}

\subsubsection{The umlaut-only class}
\label{sec:plural-umlaut-only}

Nouns in this class form the plural by umlaut alone, without an overt plural suffix. Table~\ref{tab:plural-umlaut-only} gives the forms and their vowel alternations.

\begin{table}[!ht]
    \caption{Plurals formed by umlaut alone}
    \label{tab:plural-umlaut-only}
    \small
    \begin{tabularx}{\textwidth}{l l >{\raggedright\arraybackslash}X >{\raggedright\arraybackslash}X}
        \toprule
        \rowcolor{grammarTableHead}
        \textbf{Singular} & \textbf{Plural} & \textbf{Meaning} & \textbf{Change} \\
        \midrule
        \nk{hond} & \nk{hynd} & child $\rightarrow$ children & \nk{o} $\rightarrow$ \nk{y} \\
        \nk{tadd} & \nk{tädd} & father $\rightarrow$ fathers & \nk{a} $\rightarrow$ \nk{ä} \\
        \nk{tak} & \nk{täk} & brow $\rightarrow$ brows & \nk{a} $\rightarrow$ \nk{ä} \\
        \nk{hlåd} & \nk{hlöd} & summer $\rightarrow$ summers & \nk{å} $\rightarrow$ \nk{ö} \\
        \nk{kung} & \nk{kyng} & mirror $\rightarrow$ mirrors & \nk{u} $\rightarrow$ \nk{y} \\
        \bottomrule
    \end{tabularx}
\end{table}

\subsubsection{Zero plurals}
\label{sec:plural-zero}

A small class of nouns has a zero plural, conventionally represented as \nk{-Ø}. The singular and plural have the same form, without an added ending or a change to the stem vowel. These nouns differ from the umlaut-only class, whose plural has a changed vowel. Table~\ref{tab:plural-zero} gives the zero-plural paradigms.

The class remains small and is not broadly productive.

\begin{table}[!ht]
    \caption{Zero plurals}
    \label{tab:plural-zero}
    \small
    \begin{tabularx}{\textwidth}{l l l >{\raggedright\arraybackslash}X}
        \toprule
        \rowcolor{grammarTableHead}
        \textbf{Singular} & \textbf{Plural} & \textbf{Gender} & \textbf{Meaning} \\
        \midrule
        \nk{dyr} & \nk{dyr} & Masculine & bear \\
        \nk{svern} & \nk{svern} & Masculine & red deer, stag \\
        \nk{skåv} & \nk{skåv} & Neuter & reef, submerged rock \\
        \nk{hærn} & \nk{hærn} & Neuter & antler \\
        \nk{töng} & \nk{töng} & Masculine & salmon \\
        \nk{våst} & \nk{våst} & Neuter & rocky field, boulder field \\
        \nk{norg} & \nk{norg} & Neuter & rocky islet \\
        \nk{hålv} & \nk{hålv} & Neuter & deep river pool \\
        \bottomrule
    \end{tabularx}
\end{table}

Zero plurals remain unchanged where a plural noun is required, as in \nikp{der dyr}{the bears}, beside singular \nikp{en dyr}{the bear}. This lexical pattern is distinct from the general use of the singular after a numeral or quantifier: \nk{dyr} has an unchanged plural even without a numeral, whereas \nk{afte} has the plural \nk{aften} but occurs in the singular in \nikp{ty afte}{two books}.

\section{The case system}
\label{sec:nominal-case}

Nouns in Niklorsk distinguish two cases: the nominative and the genitive. Definite articles have separate dative forms, and personal pronouns also distinguish accusative and dative forms (see Table~\ref{tab:personal-pronouns}). Recipient marking and object order are described in \S\ref{sec:ditransitive-verbs}.

The genitive case is used to indicate possession, as in the following examples:

\pex
\begingl
\gla Ga malmar at Martint brug i kvett.//
\glb \First\Pl{} go.\Pst{} to Martin.\Gen{} house in previous:day//
\glft `We went to Martin's house yesterday.'//
\endgl
\xe

The nominative case is the unmarked form of the noun and is used everywhere the genitive case is not. The genitive case is formed in the singular by adding the suffix \nk{-gt} to the nominative form of a noun ending in a vowel, \nk{-t} to a word ending in any of the consonants \nk{m}, \nk{n}, \nk{ng}, \nk{r}, \nk{l}, \nk{g}, \nk{v}, \nk{f}, \nk{sk}, or \nk{s}, and \nk{-öt} for all other nouns. In the plural, the suffix \nk{-es} is added to the genitive singular form of the noun if that form takes the suffix \nk{-t} and \nk{-gt}, and \nk{-s} if that form takes the suffix \nk{-öt}.

\begin{itemize}
    \item \nikp{en blivt avto}{the man's car} $\rightarrow$ \nikp{der blivtes avto}{the men's car}
    \item \nikp{der skällt bagen}{the fish's bones} $\rightarrow$ \nikp{der skälltes bagen}{the fishes' bones}
\end{itemize}

If the noun is part of a noun phrase, the genitive suffix is added to the last word in the phrase, which may not necessarily be the head noun.

The genitive is a phrasal enclitic: it marks the whole possessor phrase and attaches phonologically to its final word. The unstressed vowel of \nk{-öt} is reduced to [\ipa{ə}] and is omitted where the resulting consonant cluster is licensed. The consonantal genitive marker remains audible. Its pronunciation depends on the final word of the phrase, while its grammatical scope extends over the whole possessor.

The enclitic has no independent lexical accent. Its attachment does not add an accent domain to the final word of the possessor phrase.

The vowel in plural genitive \nk{-es} is likewise weak and subject to cluster-conditioned syncope. Vowel loss preserves the consonantal number and genitive marking and does not change the written distribution of \nk{-t}, \nk{-gt}, \nk{-öt}, and their plural forms.

\pex
\a \nikp{en blivt brug}{the man's house}
\a \nikp{en svorade blivt brug}{the sleeping man's house}
\a \nikp{en bliv i bart brug}{the house of the man at the bar}
\xe

The genitive forms always appear before the noun they modify. Alternatively, the noun marked in the genitive case can follow the head noun but must be preceded by the preposition \nikp{sla}{of, from}. The genitive case ending is retained despite the inversion. One has, for example, \nikp{en blivt brug}{the man's house} and \nikp{de brug slär blivt}{the house of the man}.

The use of the bare genitive and the \nk{sla} construction to indicate possession (and similar uses of the genitive case) is in general interchangeable. Nonetheless, some general tendencies can be observed. The bare genitive is preferred when the possessor is animate, especially when the possessor is a person. Where such possession refers to `true' possession, i.e., ownership etc., the bare genitive is required. Where such possession refers to an abstract concept, e.g., a quality or attribute of the animate possessor, the \nk{sla} construction is preferred if the emphasis is on the quality or attribute, while the bare genitive is preferred if the emphasis is on the possessor.

\pex
\a 
\begingl
\gla Andärst mobil sögaren i kvett.//
\glb Andärs.\Gen{} mobile lose.\Pst{}.\Pass{} in previous:day//
\glft `Andärs' mobile was lost yesterday.'//
\endgl
\a
\begingl
\gla Andärst hvær destarar ægn.//
\glb Andärs.\Gen{} charm surprise.\Pst{} \First{}\Sg{}.\Acc{}//
\glft `Andärs' charm surprised me.'//
\endgl
\a
\begingl
\gla En hvær sla Andärst destarar ægn.//
\glb the.\M{} charm of Andärs.\Gen{} surprise.\Pst{} \First{}\Sg{}.\Acc{}//
\glft `The charm of Andärs surprised me.'//
\endgl
\xe

\section{Personal pronouns}
\label{sec:personal-pronouns}

As in English, Niklorsk pronouns replace nouns or noun phrases. Unlike nouns, however, pronouns retain case marking; in addition to the unmarked nominative form and the genitive form, pronouns have distinct accusative and dative forms, used when the pronoun is the direct or indirect object of a verb, respectively, among other things. The declension of personal pronouns is shown in Table~\ref{tab:personal-pronouns}.

\begin{table}[!ht]
    \caption{Personal pronouns in Niklorsk}
    \label{tab:personal-pronouns}
    \small
    \begin{tblr}
        {
            colspec={X[2]XXXX}
        }
        Person & English & Base form & Accusative & Dative \\
        First person, sg. & `I' & \nk{ag} & \nk{ægn} & \nk{ægnar} \\
        Second person, sg. & `you' & \nk{meg} & \nk{mæg} & \nk{mæg} \\
        Third person, masc., sg. & `he' & \nk{nar} & \nk{när} & \nk{när} \\
        Third person, fem., sg. & `she' & \nk{hon} & \nk{hona} & \nk{hons} \\
        Third person, nt., sg. & `it' & \nk{jö} & \nk{ja} & \nk{järs} \\
        First person, pl. & `we' & \nk{ga} & \nk{gor} & \nk{gæ} \\
        Second person, pl. & `you' & \nk{it} & \nk{it} & \nk{itte} \\
        Third person, pl. & `they' & \nk{nör} & \nk{nörn} & \nk{nörne} \\
    \end{tblr}
\end{table}

The third person pronoun form used depends on the grammatical gender of the noun it refers to or replaces. Thus \nk{nar} is used to refer to a masculine noun, \nk{hon} to a feminine noun, and \nk{jö} to a neuter noun, unlike English, for example, where `it' is used for all inanimate objects. The third person plural pronoun \nk{nör} is used to refer to a group of people or things, regardless of the grammatical gender.

\subsection{Strong and weak pronouns}
\label{sec:pronoun-reduction}

Personal pronouns have strong pronunciations in citation, correction, and contrastive focus. Unstressed pronouns have shorter vowels and reduced vowel qualities, as shown in Table~\ref{tab:weak-pronouns}. These weak forms attach prosodically to an adjacent word without changing their syntactic position or case function.

Weak pronouns have no independent lexical accent. An independent strong pronoun can realise an accent at primary stress when its rhyme is eligible; full vowel quality alone does not establish tonal eligibility (Chapter~\ref{chap:phonology}, \S\ref{sec:accent-licensing}).

\begin{table}[!ht]
    \caption{Weak forms of personal pronouns}
    \label{tab:weak-pronouns}
    \small
    \begin{tabularx}{\textwidth}{l X l}
        \toprule
        \rowcolor{grammarTableHead}
        \textbf{Pronoun} & \textbf{Function} & \textbf{Weak pronunciation} \\
        \midrule
        \nk{ag} & First person singular nominative & [\ipa{ɐ}] \\
        \nk{meg} & Second person singular nominative & [\ipa{mə}] \\
        \nk{nar} & Third person masculine singular nominative & [\ipa{nɐ}] \\
        \nk{när} & Third person masculine singular accusative or dative & [\ipa{nɛ}] \\
        \nk{hon} & Third person feminine singular nominative & [\ipa{hən}] \\
        \nk{hona} & Third person feminine singular accusative & [\ipa{hənɐ}] \\
        \nk{jö} & Third person neuter singular nominative & [\ipa{jə}] \\
        \nk{ja} & Third person neuter singular accusative & [\ipa{jɐ}] \\
        \bottomrule
    \end{tabularx}
\end{table}

In rapid speech, \nk{ag} is further reduced to [\ipa{ə}], and weak \nk{hon} loses initial /h/, giving [\ipa{ən}]. The weak form of \nk{nör} is shortened but retains a rounded vowel. Reduction does not require every pronoun vowel to merge with schwa: \nk{nar} and \nk{när}, and \nk{jö} and \nk{ja}, retain distinct weak forms. The case paradigm in Table~\ref{tab:personal-pronouns} applies in both strong and weak positions.

\pex
\begingl
\gla Martin vor dänn brug. Nar hjekkatt av Kari.//
\glb Martin be in house. He talk.\Prog{} with Kari//
\glft `Martin is in the house. He is talking with Kari.'//
\endgl
\xe

\pex
\begingl
\gla Eg tukk ferten; jö vædaren fänt i Martint brord.//
\glb the.\Nt{} box missing.\Pass{} it see.\Pst{}.\Pass{} last in Martin.\Gen{} room.//
\glft `The box is missing; it was last seen in Martin's room.'//
\endgl
\xe


\subsection{The reflexive pronoun ham}

The reflexive pronoun \nk{ham} is used to refer back to the subject of the sentence

\section{Demonstrative pronouns}

Niklorsk has a three-way distinction in demonstrative pronouns: \nk{man} (this) is used to refer to something close to the speaker (proximal), \nk{mår} (that) is used to refer to something close to the listener (medial), and \nk{med} (that over there) is used to refer to something far from both the speaker and the listener (distal). These demonstratives may be used on their own or followed by a noun, in which case the demonstrative agrees with the noun in gender and number. Note that the forms given above correspond to the neuter singular forms of the demonstratives; nevertheless, the plural forms may also be used by itself if the implicit noun is plural.

\begin{table}[!ht]
    \caption{Declension of the demonstrative pronouns}
    \label{tab:dem-pronouns}
    \small
    \begin{tblr}{
        colspec={XXXXX}
    }
        & Masculine & Feminine & Neuter & Plural \\
        Proximal & \nk{man} & \nk{man} & \nk{man} & \nk{mannar}\\
        Medial & \nk{mörre} & \nk{mår} & \nk{mår} & \nk{mär} \\
        Distal & \nk{midda} & \nk{medder} & \nk{med} & \nk{mor}
    \end{tblr}
\end{table}

\pex
\begingl
\gla Man vase är gimlagnär i Malmö.//
\glb this.\F{} vase be.\Pst{} buy.\PstP{}.\F{} in Malmö//
\glft `This vase was bought in Malmö.' or `(I) bought this vase in Malmö.'//
\endgl
\xe

\pex
\begingl
\gla Med ö likke kort, gävs mannar ö någninna.//
\glb that.\Nt{} be.\Pres{} quite old but this.\Pl{} be.\Pres{} new.\Comp{}.\Pl{}//
\glft `That one over there is quite old, but these ones are newer.'//
\endgl
\xe

\nk{Med} (or its plural form \nk{mor}) is also often used where English might use a third person pronoun, especially when distinguish between two or more referrents. This is true also in idiomatic expressions or proverbs referring to an unspecified person or thing.


\nk{Med} is also used as an expletive pronoun, corresponding to English `it,' serving as the overt subject of an impersonal verb. In this usage, \nk{med} is undeclined.

\pex
\begingl
\gla Udbär morg med åd naups.//
\glb sometimes rain.\Pres{} it after snow.\Inf{}//
\glft `It sometimes rains after snow.'//
\endgl
\xe

\pex
\begingl
\gla Med ö skrig begns eg Andärst roghjalt.//
\glb it be.\Pres{} difficult read.\Inf{} the.\Nt{} Andärs.\Gen{} handwriting//
\glft `It is difficult to read Andärs' handwriting.'//
\endgl
\xe

If the predicate is an adjective, however, the pronoun \nk{med} may be omitted and the adjective fronted. Thus \nk{Med ö skrig begns eg Andärst roghjalt} may be rephrased as

\pex
\begingl
\gla Skrig ö begns eg Andärst roghjalt.//
\glb difficult be.\Pres{} read.\Inf{} the.\Nt{} Andärs.\Gen{} handwriting//
\glft `It is difficult to read Andärs' handwriting.'//
\endgl
\xe


The definite nominalized form \nikp{eg blä}{the same} may be used as a demonstrative, but only nominally, to refer to something previously mentioned. This usage is often emphatic and is most often found in legal writing or in other formal contexts. Nevertheless, the adjective \nikp{blä}{same} may be used attributively, without the same connotation as \nk{eg blä}.

\pex
\begingl
\gla Matti vor i Malmö örs eg blä mykatt bram imed.//
\glb Matti be.\Pres{} in Malmö and the.\Nt{} same make.\Prog{} study there//
\glft `Matti is in Malmö and he is studying there.'//
\endgl
\xe

\pex
\begingl
\gla Matti ö de blär person, käll vædar ag i Malmö.//
\glb Matti be.\Pres{} the.\F{} same.\F{} person who see.\Pst{} \First{}\Sg{} in Malmö//
\glft `Matti is the same person whom I saw in Malmö.'//
\endgl
\xe



\section{Indefinite pronouns}

\subsection{Declension}

Indefinite pronouns may be used nominally or attributively, like demonstrative pronouns.

\nk{Læse} and \nk{sonne} correspond to English `every' and `each,' respectively. The latter takes a partitive sense, i.e., it emphasizes the individual members of a group. In some contexts, however, they might be used interchangeably. When used nominally, their neuter forms \nk{læse} and \nk{sonne} are used, and they are usually translated as `everyone,' `everything,' `each one,' etc.

The form \nk{læse} is used in masculine, neuter, and plural agreement. The feminine singular form is \nk{læser}.

\begin{table}[!ht]
    \caption{Declension of indefinite pronouns}
    \label{tab:indef-pronouns}
    \small
    \begin{tblr}{
        colspec={XXXXX}
    }
       
        & Masculine & Feminine & Neuter & Plural \\
        \nikp{elge}{any} & elge & elger & elge & elge \\
        \nikp{ska}{all} & ska & ska & ska & ska \\
        \nikp{læse}{every} & \nk{læse} & \nk{læser} & \nk{læse} & \nk{læse} \\
        \nikp{sonne}{each} & sonne & sonne & sonne & sonne \\
        \nikp{väl}{some} & väl & väler & väl & väla \\
        \nikp{smä}{no} & smä & smär & smä & små \\
    \end{tblr}
\end{table}

\pex
\begingl
\gla Læse lod de ngell til målas i finns.//
\glb every.\Nt{} have.\Pres{} the right for live.\Inf{} in peace//
\glft `Everyone has the right to live in peace.'//
\endgl
\xe

\pex
\begingl
\gla Læse dägn begns man afte.//
\glb every.\Nt{} must.\Pres{} read.\Inf{} this book//
\glft `Everyone must read this book.'//
\endgl
\xe

\pex
\begingl
\gla Nar begnar sonne i afte.//
\glb he read.\Pst{} each in book//
\glft `He read each (one) in the book.' That is, `He read each page of the book.'//
\endgl
\xe

The preposition \nikp{i}{in, on, at} is used with indefinite pronouns to form locative or temporal indefinite adverbial phrases, as in English `nowhere,' `everywhere,' `anytime' etc.

\pex
\begingl
\gla I smä igom Martin hvengsen dänn kampus.//
\glb in no.\Nt{} can.\Pres{} Martin find.\Inf{}.\Pass{} in campus//
\glft `Nowhere can Martin be found on campus.'//
\endgl
\xe

\section{Interrogative pronouns}
\label{sec:interrogative-pronouns}

The interrogative pronouns in Niklorsk are shown in Table~\ref{tab:interr-pronouns} below. In general, pronouns are indeclinable with the exception of \nikp{käll}{who} and \nikp{kan}{what}, which has the plural forms \nikp{käller} and \nikp{känner}, respectively.

\begin{table}[!ht]
    \caption{Interrogative pronouns in Niklorsk}
    \label{tab:interr-pronouns}
    \small
    \begin{tblr}{
        colspec={XX},
        row{1}={font=\normalfont,bg=white},
        hline{2}={0pt},
    }
    \nikp{käll}{who} & \nikp{kan}{what} \\
    \nikp{käller}{who (pl.)} & \nikp{känner}{what (pl.)} \\
    \nikp{käg}{where} & \nikp{kamm}{when} \\
\end{tblr}
\end{table}

\enlargethispage{20bp}
\section{Syntax of noun phrases}
